\chapter{Problémafelvetés és motiváció}

\section{Miért nem elég egy hagyományos statisztikai oldal?}

A legtöbb, játékosoknak szánt statisztikai felület egyetlen mérkőzésre, egyetlen rangsorra vagy néhány alapmutatóra szűkíti le az értelmezést. Ez önmagában hasznos lehet, de nem válaszolja meg azt a kérdést, amely a projekt egészét végigkísérte: vajon a mérkőzések mögött kirajzolódik-e egy tartós társas szerkezet, és ha igen, ez a szerkezet hogyan kapcsolódik a teljesítményhez?

A \emph{Premade Graph} két egymástól nem független problémával szembesül. Az egyik adatmérnöki jellegű: a nyers Riot API válaszokból valahogy olyan adatkorpuszt kell felépíteni, ahol a játékosok azonosítása, a meccsek tárolása és a gráfprojekció újrafuttatható módon működik -- nem csak egyszer, hanem bővíthető és ellenőrizhető formában. A másik inkább kutatási: mit lehet egyáltalán állítani ezekből az adatokból, ami nem egyszerűen látványos, hanem valóban alátámasztható is.

Ez az oka annak, hogy a dolgozat fókusza tudatosan nem egyetlen UI-oldal vagy egyetlen algoritmus. A projekt értéke éppen abban áll, hogy az adatgyűjtés, a tisztítás, a pontszámítás, a gráfépítés, a klaszterezés, a futásidejű analitika és a vizualizáció ugyanannak a rendszernek az egymásra épülő rétegei.

\section{A projekt szokatlan és érdekes volta}

A szakdolgozat témája azért különösen érdekes, mert a kiinduló adathalmaz nem egyszerű játékoslista, hanem ismétlődő interakciós háló. A rendszer ugyanazt a két játékost nemcsak egyszeri ellenfelekként látja, hanem potenciális visszatérő szövetségesként, visszatérő ellenfélként, klasztertagként, hídcsúcsként vagy egy nagyobb globális hálózat részeként is.

Ez több okból ad erős szakdolgozati alapot:

\begin{itemize}
    \item a League of Legends mérkőzésadatai elég gazdagok ahhoz, hogy valódi teljesítményprofilok legyenek számíthatók;
    \item a visszatérő együttjátszás és egymás elleni előfordulás gráfelméleti szempontból értelmezhető kapcsolatokat hoz létre;
    \item a szövetségesi és ellenfél-jelleg együttese miatt a projekt nem csak súlyozott, hanem többféle kapcsolatprojekcióként is vizsgálható;
    \item a rendszer implementációs oldalon is érdekes, mert több nyelv és futtatókörnyezet közötti felelősségmegosztást kell megindokolni.
\end{itemize}

Ez a kombináció ritkább, mint elsőre látszik. Sok játékelemző projekt vagy a vizuális oldalnál áll meg, vagy csak egy-egy táblázatos mutatóig jut el. Itt viszont ugyanannak a rendszernek része a nyers adatgyűjtés, a szerepkörérzékeny score-réteg, a klaszterezés, az útvonalkeresés, az asszortativitási analitika, a Brandes-központiság és a 3D globális megjelenítés.

\section{A központi szakdolgozati állítás}

A dolgozat fő tézise az, hogy \emph{a valós mérkőzéstörténetből felépített League of Legends játékosháló egyszerre kezelhető mint asszociatív játékosgráf, mint teljesítményprofilokkal dúsított adatstruktúra, és mint rendszertechnikai szempontból komoly, többnyelvű analitikai platform}.

Ez három irányban bontható ki. A hálózati irány azt vizsgálja, hogy az ismétlődő együttjátszás és egymás elleni előfordulás mérhető szerkezeti mintázatot alkot-e -- a \texttt{flexset} adatkészleten ez leginkább asszociatív core-periphery szerkezetként mutatkozik meg, ahol a központban hídgazdag, visszatérő ally-csoportok sűrűsödnek. A metrikai irány a \texttt{feedscore} és \texttt{opscore} pontszámokra épül: ezek nem tökéletes teljesítménymutatók, de elég konzisztensek ahhoz, hogy gráfanalitikai bemenetként használhatók legyenek. A rendszertechnikai irány pedig azt indokolja, miért szükséges Python, SQLite, Node/Express, Rust és React együttese -- nem eszközpreferencia miatt, hanem mert a különböző részfeladatok eltérő természete ezt kívánja meg.

\section{Hatókör és korlátok}

A projekt viszonylag nagy rendszer, és könnyen csúszhat abba a hibába, hogy minden meglévő ötlet kész eredményként szerepeljen. A dolgozat igyekszik ezt elkerülni. Ami ténylegesen fut és ellenőrizhető, az kész állapotként szerepel. Ami megvan a repositoryban, de a jelenlegi adatpillanatképben nincs érdemi kimenete, az korábbi vagy kiegészítő modulként jelenik meg. Amit a fejlesztési tervek jövőbeli célként fogalmaznak meg, az nem szerepel kész feature-ként.

Konkrét példa: a klaszterekhez kötött ország- és régióelemzés kódja megvan, de az ellenőrzött adatbázisban a \texttt{country} mező nem tartalmazott értékelhető adatot, ezért ez a vonal kísérleti modulként van jelölve.

\section{A dolgozat szerkezete}

A dolgozat logikája a projekt fejlődését követi. Először a problémafelvetés és a szakirodalmi háttér kerül bemutatásra, majd a Riot API-ra épülő adatforrások és a korai pipeline. Ezt követően a pontozási rendszer, a gráfépítés és a régi ország-alapú elemzési ág tárgyalása következik. A középső fejezetek az architektúra evolúcióját, a több-adatkészletes modellt, az SQLite perzisztenciát és a Rust futtatókörnyezetet tárgyalják. Végül külön fejezetek mutatják be a Pathfinder Labot, a 3D bird's-eye sphere nézetet, az asszortativitási és központisági elemzéseket, valamint a strukturális egyensúly módszertani visszaminősítését, majd a validációt, a rendszerszintű értékelést és a jövőbeli fejlesztéseket.
